\documentclass{ximera}

\ifdefined\HCode
  \newcommand{\alttext}[1]{\HCode{<span class="sr-only">#1</span>}}
\else
  \newcommand{\alttext}[1]{} % Fallback for PDF view
\fi



% MATH -----------------------------------------------------------
\newcommand{\nats}{\mathbb N}
\newcommand{\ints}{\mathbb Z}
\newcommand{\rats}{\mathbb Q}
\newcommand{\reals}{\mathbb R}
\newcommand{\complex}{\mathbb C}
\newcommand{\powerset}{\mathscr P}

%-------------------------------------------------------------

\pgfplotsset{compat=1.18}
\begin{document}
\begin{abstract}
 \end{abstract}

    \title{Class Discussion Activity 18}
    
    \maketitle

\section*{Exercises}

\begin{enumerate}[label=(\arabic*)]


\item A large undamped (vertical) spring-mass system begins (at time $t=0$) at rest at equilibrium.  Then the mass is constantly forced upward by a force of magnitude $12$ Newtons.  Suppose the mass is $2$ kg and the spring stiffness is $8$ Newtons per meter.  Find the vertical displacement of the mass (in meters) at time $t=10$ seconds.  Round to the nearest hundredth of a meter.

    % 2y''+8y=12 so y''+4y=6.  y=1.5-1.5cos(2t).  So y(10)=1.5(1-cos(20))\approx 0.89


\item Which CANNOT be displacement of the mass above equilibrium for a free damped spring-mass system?

\begin{enumerate}[label=(\alph*)]

\item $y=(1+t)e^{-t/2}$ 
\item $y=2e^{-t}-e^{-10t}$
\item $y=e^{-0.2t}\cos{(\pi t)}$
\item $y=e^{-t/10}-10e^{-t/100}$ 
\item $y=\dfrac{\cos{(2\pi t)}-\sin{(2\pi t)}}{e^{t}}$ 
\item $y=\dfrac{1+t}{e^{1-t/5}}$

\end{enumerate}

% (f) with options through (f)

\newpage


\item Each curve below is the graph of the displacement above equilibrium of a mass in a spring-mass system.  Match each curve to the types of spring-mass systems below.

\begin{figure}[h]
    \centering
\alttext{Some solution curves to spring-mass systems.}
\begin{tikzpicture}[scale=1.5]
    \begin{axis}[axis equal,axis lines = middle,xlabel = \(t\),grid=both,
    grid style={line width=.1pt, draw=gray!10},
    major grid style={line width=.2pt,draw=gray!50},
    ylabel = {\(y(t)\)}]
    \addplot[domain=0:4,
    samples=120,color=blue,ultra thick]{x*sin(5*deg(x))};
    \addlegendentry{\((i)\)}        
    \addplot[domain=0:4,
    samples=120,color=blue]{2*sin(8*deg(x))};
    \addlegendentry{\((ii)\)}
    \addplot[domain=0:4,
    samples=60,color=black,dotted,thick]{3+exp(-x)*cos(8*deg(x))};
    \addlegendentry{\((iii)\)}
    \addplot[domain=0:4,
    samples=60,color=black,dashed,ultra thick]{1.5(x-0.4)*exp(-x)};
    \addlegendentry{\((iv)\)}
    \addplot[domain=0:4,
    samples=60,color=black,thick]{1.5*exp(-x)*sin(10*deg(x))};
    \addlegendentry{\((v)\)}
        \end{axis}
\end{tikzpicture}
    \caption{Some solution curves to spring-mass systems}
\end{figure}

\begin{enumerate}[label=(\alph*)]
\item A simple harmonic oscillator
\item A free underdamped spring-mass system 
\item A free critically damped spring-mass system 
\item A forced undamped spring-mass system in resonance 
\item A forced and underdamped spring-mass system 
\end{enumerate}


% Matching:  (a)-(ii), (b)-(v), (c) - (iv), (d)-(i) and (e)-(iii)

\newpage


\item A small object of mass 1 kg is attached to an elastic spring of stiffness $1$ N/m and is immersed in a viscous medium that resists the motion of the mass with 2 Newtons of force per m/s of speed.  Suppose that at a certain instant the spring is $0.25$ meters below equilibrium and moving upward at $1$ meter per second.  How many times does the mass cross through equilibrium thereafter?

\begin{enumerate}[label=(\alph*)]
    \item Zero times
    \item One time % ***
    \item Two times
    \item Three times
    \item Infinitely-many times
\end{enumerate}

% (b) with options through (e)

\item The gun of a U.S. M60 tank is attached to spring-mass system with a $100$ kg mass and a dashpot damper.  The spring stiffness is engineered to equal $\dfrac{c^2}{400}$ where $c$ is the damping constant of the dashpot damper.  Assume that when the gun is fired, the mass is at equilibrium with a velocity of $100$ m/s.  It is desired that one second later, the quantity $y^2+(y\,^{\prime})^2$ be no greater than $0.01$.  What is the minimum possible value of $c$ in $N\,s/m$?  Report the answer with 2 significant figures (e.g. $890$ for $893$).

Hint: use a numerical inequality solver (e.g. Wolfram Alpha)
    
    % 1800
    
    % y=100te^(-ct/200) ...   y(1)=100e^(-c/200) and y'(1)=[100-c/2]e^(-c/200)
    
    % (100e^(-c/200))^2+([100-c/2]e^(-c/200))^2 is at least 0.01 as c goes from 1798 to 1799 ... so 1800.


\end{enumerate}

\end{document}